\documentclass[12pt,a4paper]{report}

\usepackage[utf8]{inputenc}
\usepackage[T1]{fontenc}
\usepackage[french]{babel}
\usepackage{lmodern}
\usepackage[top=2.5cm,bottom=2.5cm,left=3cm,right=2.5cm]{geometry}
\usepackage{setspace}
\onehalfspacing
\usepackage{parskip}
\setlength{\parindent}{0pt}

\usepackage{xcolor}
\usepackage{colortbl}
\definecolor{primary}{RGB}{30,80,110}
\definecolor{secondary}{RGB}{52,120,70}
\definecolor{accent}{RGB}{180,80,30}
\definecolor{lightblue}{RGB}{220,235,245}
\definecolor{lightgreen}{RGB}{220,240,225}
\definecolor{codebg}{RGB}{248,248,252}

\usepackage{booktabs}
\usepackage{tabularx}
\usepackage{array}
\newcolumntype{L}[1]{>{\raggedright\arraybackslash}p{#1}}
\newcolumntype{C}[1]{>{\centering\arraybackslash}p{#1}}

\usepackage{listings}
\lstset{
  backgroundcolor=\color{codebg},
  basicstyle=\footnotesize\ttfamily,
  keywordstyle=\color{primary}\bfseries,
  commentstyle=\color{gray}\itshape,
  stringstyle=\color{secondary},
  numberstyle=\tiny\color{gray},
  numbers=left, numbersep=8pt,
  frame=single, rulecolor=\color{gray!30},
  breaklines=true, showstringspaces=false,
  language=Python
}

\usepackage[hidelinks]{hyperref}
\usepackage{fancyhdr}
\pagestyle{fancy}
\fancyhf{}
\fancyhead[L]{\small\nouppercase{\leftmark}}
\fancyhead[R]{\small TER M1 --- Agentic AI}
\fancyfoot[C]{\thepage}
\renewcommand{\headrulewidth}{0.4pt}

\usepackage[authoryear,round,sort]{natbib}

\usepackage{enumitem}
\usepackage{amsmath}
\usepackage{float}

% Boîtes légères sans mdframed
\newenvironment{infobox}[1]{%
  \begin{center}
  \colorbox{lightblue}{\parbox{0.95\textwidth}{%
  \medskip\textbf{\color{primary}#1}\par\smallskip}}\end{center}%
  \begin{center}\colorbox{lightblue}{\parbox{0.95\textwidth}\bgroup\smallskip
}{%
  \smallskip\egroup}\end{center}
}

\newenvironment{defbox}[1]{%
  \begin{center}\colorbox{lightgreen}{\parbox{0.95\textwidth}{%
  \smallskip\textbf{\color{secondary}#1}\par\smallskip
}{%
  \smallskip}}\end{center}
}

% ─────────────────────────────────────────────────────────────────────────────
\begin{document}

% ── Page de titre ─────────────────────────────────────────────────────────────
\begin{titlepage}
\begin{center}
\vspace*{1cm}
{\Large\bfseries Université --- Master Informatique}\\[0.3em]
{\large Intelligence Artificielle \& Data}\\[0.5em]
\rule{\textwidth}{1.5pt}\\[1.5em]
{\LARGE\bfseries\color{primary}
Architectures Logicielles pour l'IA Agentique\\[0.5em]}
{\large\bfseries Systèmes Multi-Agents, Variabilité Architecturale\\
et Prototype de Recrutement Automatisé}\\[0.5em]
\rule{\textwidth}{1.5pt}\\[1.5em]
{\large Travaux d'Étude et de Recherche (TER) --- Master 1}\\[0.4em]
{\large Année universitaire 2025--2026}\\[2cm]
\begin{tabular}{ll}
\textbf{Auteurs :} & Sofiane Dzermane\\
& [Prénom Nom]\\[1em]
\textbf{Encadrant :} & [Nom de l'encadrant]\\[1em]
\textbf{Date :} & Mai 2026\\
\end{tabular}
\end{center}
\end{titlepage}

% ── Résumé ────────────────────────────────────────────────────────────────────
\chapter*{Résumé}
\addcontentsline{toc}{chapter}{Résumé}

Ce rapport présente les travaux réalisés dans le cadre du TER de Master~1
Informatique, portant sur l'étude et la mise en \oe uvre d'architectures logicielles
pour les systèmes basés sur l'intelligence artificielle agentique.

La première partie constitue un état de l'art approfondi couvrant deux domaines
complémentaires : les Systèmes Multi-Agents (SMA) classiques et l'Agentic~AI
contemporaine fondée sur les modèles de langage (LLM). Une analyse comparative
rigoureuse, appuyée sur dix-sept articles scientifiques, identifie les convergences
et divergences entre ces paradigmes.

La seconde partie présente la conception et le développement d'un prototype de
\textbf{système de recrutement automatisé multi-agents}. Ce prototype implémente
huit agents spécialisés orchestrés par LangGraph, une mémoire contextuelle
vectorielle inter-runs (ChromaDB), un mécanisme de validation humaine
(Human-in-the-Loop), deux modes d'architecture déployables, une API~REST avec
streaming SSE, une interface graphique NiceGUI et une suite de 91 tests automatisés.

\textbf{Mots-clés :} Systèmes Multi-Agents, Agentic AI, LangGraph, RAG, LLM,
Feature Model, Microservices, Human-in-the-Loop, SE4AI, AI4SE.

\chapter*{Abstract}
\addcontentsline{toc}{chapter}{Abstract}

This report presents the work carried out as part of the Master~1 Computer Science
research project (TER), focusing on the study and implementation of software
architectures for agentic AI-based systems.

The first part constitutes a thorough state of the art covering two complementary
domains: classical Multi-Agent Systems (MAS) and contemporary Agentic~AI based on
Large Language Models. The second part presents the design and development of an
automated multi-agent recruitment system with eight agents, LangGraph orchestration,
cross-run vector memory (ChromaDB), HITL validation, microservices Docker, and
91 automated tests.

\textbf{Keywords:} Multi-Agent Systems, Agentic AI, LangGraph, RAG, Feature Model,
Microservices, Human-in-the-Loop, SE4AI.

\tableofcontents
\listoffigures
\listoftables

% =============================================================================
\chapter{Introduction}
\label{chap:intro}
% =============================================================================

\section{Contexte et motivation}

L'essor des modèles de langage de grande taille (LLM) depuis 2022 a profondément
transformé la façon dont les systèmes logiciels intègrent l'intelligence artificielle.
Au-delà de la simple génération de texte, ces modèles servent aujourd'hui de
\textit{couche cognitive} dans des systèmes distribués capables de planifier, percevoir
leur environnement et orchestrer des outils : les \textbf{systèmes agentiques}.

Ce mouvement s'inscrit dans une dialectique entre deux communautés :

\begin{itemize}
  \item Le \textbf{génie logiciel} (GL) : concevoir des systèmes fiables, maintenables
        et performants via des principes d'architecture et des pratiques d'ingénierie.
  \item L'\textbf{intelligence artificielle} (IA) : produire des capacités
        d'apprentissage et de décision autonome, du machine learning aux LLM.
\end{itemize}

Ces deux mondes interagissent selon deux axes :

\begin{description}
  \item[AI4SE] (\textit{Artificial Intelligence for Software Engineering}) :
        l'IA améliore les activités d'ingénierie logicielle.
  \item[SE4AI] (\textit{Software Engineering for AI}) :
        le génie logiciel apporte méthodes et architectures pour industrialiser les
        systèmes IA.
\end{description}

\section{Objectifs du TER}

\begin{enumerate}
  \item Étudier la \textbf{variabilité architecturale} des systèmes logiciels basés
        sur l'IA et la représenter sous la forme d'un feature model.
  \item Identifier des \textbf{critères de choix} orientant la sélection d'un style
        architectural.
  \item Implémenter une \textbf{architecture de référence Agentic AI} fondée sur les
        microservices et son écosystème (conteneurisation, API Gateway, RAG,
        observabilité).
\end{enumerate}

\section{Démarche suivie}

\begin{enumerate}
  \item \textbf{Phase de recherche bibliographique} : étude de dix-sept articles
        scientifiques sur les SMA classiques et l'Agentic AI.
  \item \textbf{Phase de choix du scénario} : sélection du recrutement automatisé
        parmi plusieurs candidats.
  \item \textbf{Phase de conception et développement} : architecture, implémentation,
        tests et évaluation du prototype.
\end{enumerate}

\section{Organisation du rapport}

Le chapitre~\ref{chap:soa} présente l'état de l'art.
Le chapitre~\ref{chap:sujet} justifie le choix du scénario.
Le chapitre~\ref{chap:conception} détaille la conception.
Le chapitre~\ref{chap:impl} décrit l'implémentation.
Le chapitre~\ref{chap:eval} présente l'évaluation.
Le chapitre~\ref{chap:conclusion} conclut.

% =============================================================================
\chapter{État de l'art}
\label{chap:soa}
% =============================================================================

\section{Systèmes Multi-Agents classiques}
\label{sec:sma}

\subsection{Définition et fondements}

Selon \citet{dorri2018}, un Système Multi-Agents (SMA) est un ensemble d'agents
autonomes interagissant dans un environnement partagé pour atteindre des objectifs
individuels ou collectifs. \citet{maldonado2024} identifient cinq composants clés :
\textbf{Agents}, \textbf{Environnement}, \textbf{Interactions}, \textbf{Organisation}
et \textbf{Objectifs}.

\medskip
\colorbox{lightgreen}{\parbox{0.95\textwidth}{
\smallskip
\textbf{\color{secondary}Propriétés fondamentales d'un agent}\\
Tout agent possède : \textbf{Autonomie}, \textbf{Réactivité}, \textbf{Proactivité}
et \textbf{Sociabilité} \citep{maldonado2024}.
\smallskip
}}
\medskip

\subsection{Architectures organisationnelles}

\citet{maldonado2024} distinguent trois styles majeurs :

\begin{itemize}
  \item \textbf{Centralisée} : un agent maître coordonne les subordonnés. Avantage :
        cohérence globale. Inconvénient : point unique de défaillance.
  \item \textbf{Décentralisée} : la coordination émerge des interactions locales.
        Avantage : robustesse. Inconvénient : comportement global difficile à prédire.
  \item \textbf{Hiérarchique} : multi-niveaux, compromis le plus répandu.
\end{itemize}

\subsection{Coordination et protocoles formels}

\citet{generalsurvey2020} formalisent les mécanismes d'interaction :
Contract Net Protocol, théorie des jeux, enchères, consensus distribué et
ACL (\textit{Agent Communication Language}) \citep{dorri2018}.

\subsection{Apports du Machine Learning : le MARL}

\citet{cmupaper2021} établissent un pont via le Multi-Agent Reinforcement Learning
(MARL). Les défis principaux sont la non-stationnarité, l'instabilité de convergence,
et l'équilibre coopération vs compétition.

\subsection{Émergence et intelligence collective}

\citet{visualsurvey2019} mettent en lumière l'émergence : des comportements globaux
complexes naissent d'interactions locales simples, dans des domaines variés : trafic,
finance, épidémiologie, smart cities \citep{wjarrpaper2024}.

\section{Agentic AI contemporaine}
\label{sec:agentic}

\subsection{Clarification conceptuelle}

L'Agentic AI ne se réduit ni à un agent LLM augmenté, ni aux SMA classiques.
Plusieurs niveaux émergent \citep{sapkota2024,schneider2025,pati2025} :

\begin{description}
  \item[AI Agents] \citep{sapkota2024} : systèmes modulaires activés par des LLM,
        spécialisés via prompt engineering et intégration d'outils.
  \item[Agentic AI] \citep{schneider2025,pati2025} : paradigme orienté objectifs,
        planification autonome, mémoire persistante et adaptation dynamique.
  \item[Communautés agentiques] \citep{milosevic2026} : écosystèmes hybrides
        humains-IA avec gouvernance formelle.
\end{description}

La trajectoire évolutive identifiée est :
\[ \text{GenAI} \to \text{AI Agents} \to \text{Agentic AI} \to
   \text{Communautés agentiques} \to \text{Écosystèmes hybrides gouvernés} \]

\subsection{Double paradigme : symbolique vs neuronal}

\citet{abouali2026} proposent un cadre analytique :

\begin{itemize}
  \item \textbf{Paradigme symbolique} : planification déterministe, vérification
        formelle, forte prédictibilité.
  \item \textbf{Paradigme neuronal} : orchestration LLM, génération probabiliste,
        adaptabilité élevée.
\end{itemize}

Les travaux convergent vers une trajectoire \textbf{neuro-symbolique hybride}.

\subsection{Architecture générale}

\citet{wavestone2024}, \citet{alvapandey2024} et \citet{vaidhyanathan2026} décrivent
une architecture modulaire : moteur de raisonnement (LLM), mémoire persistante
(bases vectorielles), planification multi-étapes, sélection dynamique d'outils,
boucle itérative analyse-plan-action-évaluation.

\subsection{Enjeux transversaux}

\citet{hughes2025} et \citet{patel2025} identifient : hallucination, sécurité
(prompt injection, fuite de données), responsabilité, consommation énergétique,
alignement éthique.

\section{Analyse comparative SMA / Agentic AI}
\label{sec:comparatif}

\subsection{Convergences}

Les deux paradigmes partagent \citep{analyse_comparative} :
autonomie des agents, poursuite active d'objectifs, coordination multi-agents,
mémoire persistante, domaines d'application communs.

\subsection{Divergences structurelles}

\begin{table}[H]
\centering
\caption{Comparaison SMA classiques vs Agentic AI}
\label{tab:comparatif}
\begin{tabularx}{\textwidth}{L{3.2cm}L{5cm}L{5cm}}
\toprule
\textbf{Critère} & \textbf{SMA classiques} & \textbf{Agentic AI} \\
\midrule
Raisonnement & Logique, déterministe & Probabiliste (LLM) \\
\addlinespace
Coordination & Protocoles formels (ACL, enchères) & Prompt engineering, tool-calling \\
\addlinespace
Apprentissage & Optionnel (MARL) & Central (LLM pré-entraîné) \\
\addlinespace
Maturité & Décennies, bases mathématiques & Écosystème fragmenté \\
\addlinespace
Explicabilité & Règles traçables & Déficit inhérent aux LLM \\
\bottomrule
\end{tabularx}
\end{table}

\subsection{Apports des SMA à l'Agentic AI}

Les SMA constituent un réservoir de solutions formelles \citep{analyse_comparative} :
protocoles de coordination, théorie des jeux \citep{generalsurvey2020},
modèle BDI \citep{maldonado2024}, gouvernance formelle \citep{milosevic2026},
MARL \citep{cmupaper2021}.

\section{Comparaison des frameworks Agentic AI}
\label{sec:frameworks}

\begin{table}[H]
\centering
\caption{Comparaison des frameworks Agentic AI \citep{vaidhyanathan2026}}
\label{tab:frameworks}
\begin{tabularx}{\textwidth}{L{2.3cm}L{2.5cm}C{1.8cm}C{1.5cm}L{4cm}}
\toprule
\textbf{Framework} & \textbf{Coordination} & \textbf{État partagé} & \textbf{HITL} & \textbf{Points forts} \\
\midrule
\textbf{LangGraph} & Graphe orienté & Riche (TypedDict) & Natif & Contrôle fin, traçabilité \\
CrewAI & Rôles/délégation & Partiel & Limité & Prise en main rapide \\
AutoGen & Conversation & Limité & Oui & Flexibilité, multi-LLM \\
OpenAI Swarm & Handoff & Absent & Non & Légèreté, simplicité \\
MetaGPT & SOP rigides & Oui & Partiel & Processus logiciel \\
\bottomrule
\end{tabularx}
\end{table}

\textbf{Justification du choix de LangGraph.} Le graphe orienté est la structure la
plus expressive pour modéliser un SMA. L'état TypedDict implémente nativement le
pattern Blackboard. Le HITL et le checkpointing sont natifs. LangGraph est
indépendant du fournisseur LLM, évitant le vendor lock-in \citep{vaidhyanathan2026}.

% =============================================================================
\chapter{Choix du scénario applicatif}
\label{chap:sujet}
% =============================================================================

\section{Critères de sélection}

\begin{enumerate}
  \item \textbf{Réalisme applicatif} : cas d'usage professionnel crédible.
  \item \textbf{Complexité multi-agents} : plusieurs agents avec rôles distincts.
  \item \textbf{Justification du pipeline} : chaque étape = un agent (SRP).
  \item \textbf{Pertinence du RAG} : mémoire inter-runs à valeur démontrable.
  \item \textbf{Faisabilité} : APIs gratuites, prototype académique.
\end{enumerate}

\section{Candidats étudiés}

\begin{table}[H]
\centering
\caption{Scénarios applicatifs évalués}
\label{tab:scenarios}
\begin{tabularx}{\textwidth}{L{3.2cm}L{3.8cm}L{4cm}C{1.8cm}}
\toprule
\textbf{Scénario} & \textbf{Points forts} & \textbf{Limites} & \textbf{Retenu} \\
\midrule
Comparateur de prix & Scraping simple & Peu de valeur LLM & Non \\
Assistant juridique & RAG pertinent & Données sensibles & Non \\
Veille concurrentielle & Multi-sources & Périmètre flou & Non \\
\textbf{Recrutement automatisé} & Pipeline riche, HITL naturel, RAG utile & LinkedIn bloqué & \textbf{Oui} \\
\bottomrule
\end{tabularx}
\end{table}

\section{Justification du recrutement automatisé}

\paragraph{Pipeline naturellement multi-agents.}
Le recrutement se décompose en étapes distinctes : analyse de poste, stratégie,
collecte, filtrage, évaluation, vérification, rédaction, persistance. Chaque étape
= un agent à responsabilité unique (principe SRP appliqué aux agents).

\paragraph{Human-in-the-Loop pertinent.}
La décision de contacter un candidat est à forte valeur métier : elle justifie une
validation humaine explicite avant action, exactement le pattern HITL à démontrer.

\paragraph{RAG démontrable.}
Les profils invalides récurrents entre runs justifient la blacklist inter-runs.
L'économie de scraping et d'appels LLM est mesurable.

% =============================================================================
\chapter{Conception et architecture}
\label{chap:conception}
% =============================================================================

\section{Vue d'ensemble du système}

Le système comprend huit agents orchestrés par LangGraph, une mémoire vectorielle
(ChromaDB), une API REST (FastAPI), une interface (NiceGUI) et deux modes de
déploiement.

\begin{figure}[H]
\centering
\begin{tabularx}{\textwidth}{X}
\toprule
\textbf{Pipeline multi-agents (séquentiel puis parallèle)} \\
\midrule
\textbf{START} $\to$ A1 Orchestrateur $\to$ A2 Analyste $\to$ A3a Stratège \\
$\to$ A3b Collecteur \textit{(DDG + GitHub)} $\to$ A3c Filtre \textit{(blacklist RAG)} $\to$ A6 Déduplicateur \\
$\to$ \textit{Send()$\times$N} $\to$ \textbf{A4 Évaluateur} \textit{(Groq LLM, cache RAG)} \\
$\to$ \textit{Send()$\times$N} $\to$ \textbf{A5 Vérificateur} \textit{(Groq LLM + guards)} \\
$\to$ Injection RAG \textit{(ChromaDB, cosine $\geq 0.75$)} $\to$ A7 Recruteur $\to$ A8 Persistance \textit{(ChromaDB)} $\to$ \textbf{END} \\
\bottomrule
\end{tabularx}
\caption{Architecture du pipeline multi-agents}
\label{fig:pipeline}
\end{figure}

\section{Feature model : variabilité architecturale}
\label{sec:featuremodel}

\begin{figure}[H]
\centering
\begin{tabularx}{\textwidth}{L{3cm}C{1.2cm}L{8.5cm}}
\toprule
\textbf{Feature} & \textbf{Type} & \textbf{Variantes} \\
\midrule
\rowcolor{lightblue}
\textbf{SMA Recrutement} & racine & --- \\
Mode d'exécution & XOR & \textbf{Local} (monolithique, LangGraph) $|$ \textbf{Microservices} (Docker) \\
Fournisseur LLM & mand. & \textbf{Groq} (\texttt{llama-3.3-70b-versatile}) \\
HITL & XOR & \textbf{Activé} (interruption avant A7) $|$ \textbf{Désactivé} (full auto) \\
Interface & OR & \textbf{CLI} $|$ \textbf{API REST} $|$ \textbf{UI Web} \\
Mémoire RAG & AND & \textbf{Blacklist} inter-runs $+$ \textbf{Cache} de scores \\
\bottomrule
\multicolumn{3}{l}{\scriptsize mand.~= mandatoire \quad XOR~= alternatif exclusif \quad OR~= un ou plusieurs \quad AND~= tous}
\end{tabularx}
\caption{Feature model du système de recrutement multi-agents}
\label{fig:featuremodel}
\end{figure}

\subsection{Décisions architecturales justifiées}

\paragraph{Mode d'exécution (XOR).}
\begin{itemize}
  \item \textbf{Local} : tous les agents dans un processus Python (LangGraph).
        Avantage : performance optimale, pas d'overhead réseau.
  \item \textbf{Microservices} : chaque agent est un service HTTP Docker.
        Avantage : scaling horizontal, déployabilité cloud.
        Inconvénient : overhead réseau (+30--50\,\%).
\end{itemize}

\paragraph{LangGraph comme bus de messages.}
Nous n'utilisons pas Kafka ou RabbitMQ. Le graphe LangGraph joue ce rôle : chaque
arête est un canal typé, le state est le message. Ce choix simplifie le déploiement
tout en conservant les propriétés d'un bus (découplage, traçabilité, replay).

\section{Patterns SMA implémentés}

\begin{table}[H]
\centering
\caption{Patterns SMA implémentés dans le prototype}
\label{tab:patterns}
\begin{tabularx}{\textwidth}{L{3cm}L{4cm}L{6cm}}
\toprule
\textbf{Pattern} & \textbf{Implémentation} & \textbf{Justification} \\
\midrule
Superviseur & A1 via flux du graphe & Topologie étoile, délégation formelle \\
Blackboard & \texttt{GraphState} TypedDict & État partagé, droits séparés par agent \\
Send/Map-Reduce & A4+A5 via \texttt{Send()} & $N$ instances parallèles \\
Validation pair-à-pair & A4 score, A5 contrôle & Séparation évaluation/vérification \\
Routage conditionnel & \texttt{route\_apres\_verification()} & Décision selon meilleur score \\
Human-in-the-Loop & \texttt{interrupt\_before=["recruteur"]} & Suspension pour validation \\
\bottomrule
\end{tabularx}
\end{table}

\section{Choix technologiques}

\begin{table}[H]
\centering
\caption{Justification des choix technologiques}
\label{tab:technos}
\begin{tabularx}{\textwidth}{L{2.5cm}L{2.5cm}L{7.5cm}}
\toprule
\textbf{Techno} & \textbf{Alternative} & \textbf{Justification} \\
\midrule
LangGraph & CrewAI, AutoGen & Graphe = SMA formel, état riche, fan-out, HITL natif \\
Groq (cloud) & Ollama (local) & Ollama crashait la machine (RAM), Groq free tier 100k tok/jour \\
ChromaDB & Pinecone, Qdrant & Local, gratuit, sans serveur \\
FastAPI & Flask & Async natif, SSE simplifié, swagger auto \\
DuckDuckGo & SerpAPI & Gratuit, sans clé API \\
LangGraph & Kafka & LangGraph est le bus : arête = canal, state = message \\
\bottomrule
\end{tabularx}
\end{table}

% =============================================================================
\chapter{Implémentation}
\label{chap:impl}
% =============================================================================

\section{Structure du projet}

\begin{description}[font=\ttfamily]
  \item[src/agents/] 8 agents + nœud injection RAG ($\sim$1\,400 lignes)
  \item[src/tools/] RAG, scraping, recherche DDG, GitHub/SO APIs
  \item[src/ui/] Interface NiceGUI (3 fichiers)
  \item[src/api.py] API Gateway FastAPI + SSE + SQLite
  \item[src/graph.py] Construction du graphe LangGraph
  \item[src/state.py] Blackboard TypedDict avec reducers
  \item[services/] 6 microservices HTTP FastAPI
  \item[tests/] 91 tests automatisés
\end{description}

\section{Le graphe LangGraph}

\[
\texttt{START} \to A1 \to A2 \to A3a \to A3b \to A3c \to A6
\xrightarrow{\text{Send()}\times N} A4 \to \texttt{reduce}
\xrightarrow{\text{Send()}\times N} A5 \to \texttt{reduce}
\to \texttt{rag} \to A7 \to A8 \to \texttt{END}
\]

Les reducers \texttt{operator.add} garantissent la cohérence du blackboard sans verrou :

\begin{lstlisting}[caption={src/state.py --- reducers pour le fan-out}]
candidats_scores:  Annotated[list[CandidatScore],  operator.add]
candidats_valides: Annotated[list[CandidatValide], operator.add]
profils_bruts:     Annotated[list[Candidat],       operator.add]
\end{lstlisting}

\section{Les huit agents}

\subsection{A1 --- Orchestrateur}
Point d'entrée. Initialise le pipeline, produit le rapport final Markdown.
Pattern Superviseur : coordination via flux du graphe.

\subsection{A2 --- Analyste de poste}
Extrait un profil JSON : hard skills, soft skills, expérience, localisation.
Enrichissement déterministe : regex pour les niveaux, stop-words pour les lieux
(``AWS'' n'est pas une ville --- bug corrigé).

\subsection{A3a --- Stratège de recherche}
Génère un plan de recherche par source (DDG, LinkedIn, GitHub, SO).
Opérateurs déterministes : \texttt{site:}, \texttt{intitle:}, \texttt{language:}.

\subsection{A3b --- Collecteur}
Exécute le plan : DDG, GitHub API officielle, Stack Overflow API.
Délai 1,5~s entre appels DDG (rate limit). Déduplication par URL.

\subsection{A3c --- Filtre anti-bruit}
Filtre algorithmique pur (sans LLM) en trois passes :
(1) pré-filtre URL (30+ agrégateurs connus),
(2) blacklist RAG inter-runs (ChromaDB metadata-only, très rapide),
(3) post-filtre mots-clés offres d'emploi.

\subsection{A6 --- Déduplicateur}
Détection par URL identique ou similarité de nom ($>85\%$,
\texttt{SequenceMatcher}). Profils fusionnés = données concaténées.

\subsection{A4 --- Évaluateur ($\times N$ parallèles)}
Scoring multicritères (0--100) : hard skills, soft skills, expérience, culture.
Cache RAG : si l'URL a déjà été évaluée pour un poste similaire (cosine $\geq 0.85$),
retourne le score sans appel LLM.

\subsection{A5 --- Vérificateur ($\times N$ parallèles)}
Validation en deux couches :
\begin{itemize}
  \item LLM : cohérence du score, détection d'incohérences.
  \item Guards déterministes : non-candidat, compétences minimales, expérience
        incompatible, upgrade auto LinkedIn \texttt{/in/} si score A4 $\geq 75$.
\end{itemize}

\subsection{Nœud Injection RAG}
Interroge ChromaDB pour les candidats valides des runs précédents sur fiches
similaires (cosine $\geq 0.75$, $\leq 30$ jours). Enrichit les résultats courants.

\subsection{A7 --- Recruteur}
Rédige des messages personnalisés. Mode absolu : score $\geq 75$.
Mode relatif : top-3 si aucun n'atteint le seuil.

\subsection{A8 --- Persistance RAG}
Stocke valides et invalides dans ChromaDB. Le statut ``invalide'' (rejet HITL)
l'emporte sur ``valide'' lors de la déduplication.

\section{Mémoire contextuelle RAG}

ChromaDB (\texttt{data/chromadb/}) avec SentenceTransformers
(\texttt{all-MiniLM-L6-v2}). Deux collections : \texttt{candidats\_evalues}
et \texttt{fiches\_poste}.

\medskip
\colorbox{lightblue}{\parbox{0.95\textwidth}{
\smallskip
\textbf{\color{primary}Propriété d'apprentissage cumulatif}\\
Le système accumule de la connaissance opérationnelle : chaque run enrichit
la mémoire et rend les suivants plus rapides et plus précis. C'est un
apprentissage par mémorisation à rendement croissant, sans modification des
poids du modèle.
\smallskip
}}
\medskip

Trois fonctions actives :
\begin{enumerate}
  \item \textbf{Blacklist (A3c)} : metadata-only, économie de scraping + LLM.
  \item \textbf{Cache (A4)} : score stocké si fiche similaire (cosine $\geq 0.85$).
  \item \textbf{Injection proactive} : candidats connus réinjectés automatiquement.
\end{enumerate}

\section{Human-in-the-Loop}

\begin{lstlisting}[caption={src/graph.py --- compilation avec HITL}]
app = graph.compile(
    checkpointer=MemorySaver(),
    interrupt_before=["recruteur"]
)
\end{lstlisting}

Trois actions via \texttt{POST /runs/\{id\}/hitl} :
\texttt{approve} (reprend), \texttt{skip} (invalide tout), \texttt{edit} (modifie
la liste en JSON).

\section{Mode microservices}

\begin{table}[H]
\centering
\caption{Services Docker}
\label{tab:microservices}
\begin{tabularx}{\textwidth}{L{3.5cm}C{1.5cm}L{8cm}}
\toprule
\textbf{Service} & \textbf{Port} & \textbf{Rôle} \\
\midrule
svc-analyste   & 8001 & A2 : analyse fiche de poste \\
svc-chercheur  & 8002 & A3 : collecte + blacklist RAG \\
svc-evaluateur & 8003 & A4 : évaluation + cache RAG \\
svc-verificateur & 8004 & A5 : vérification + guards \\
svc-recruteur  & 8005 & A7 : rédaction messages \\
svc-orchestrateur & 8006 & Orchestration HTTP \\
api-gateway    & 8000 & Point d'entrée public \\
sma-ui         & 8080 & Interface NiceGUI \\
\bottomrule
\end{tabularx}
\end{table}

\begin{table}[H]
\centering
\caption{Comparaison monolithique vs microservices}
\label{tab:tradeoffs}
\begin{tabularx}{\textwidth}{L{3cm}L{5.5cm}L{5cm}}
\toprule
\textbf{Critère} & \textbf{Monolithique (LangGraph)} & \textbf{Microservices (Docker)} \\
\midrule
Latence & Optimale (in-process) & +30--50\,\% \\
Déployabilité & Simple & Cloud-ready \\
Scaling & Vertical & Horizontal par service \\
Débogage & Traces LangGraph & Logs distribués \\
\bottomrule
\end{tabularx}
\end{table}

\section{Tests automatisés}

\begin{table}[H]
\centering
\caption{Couverture des 91 tests}
\label{tab:tests}
\begin{tabularx}{\textwidth}{L{5.5cm}C{2cm}L{5.5cm}}
\toprule
\textbf{Fichier} & \textbf{Tests} & \textbf{Couverture} \\
\midrule
\texttt{test\_graph.py} & 28 & Topologie, 16 nœuds, arêtes \\
\texttt{test\_verificateur\_rules.py} & 5 & Guards déterministes A5 \\
\texttt{test\_stratege\_queries.py} & 15 & Requêtes A3a \\
\texttt{test\_api.py} & 43 & Endpoints, HITL, SQLite \\
\midrule
\textbf{Total} & \textbf{91} & \textbf{91/91 passent} \\
\bottomrule
\end{tabularx}
\end{table}

% =============================================================================
\chapter{Évaluation}
\label{chap:eval}
% =============================================================================

\section{Run de référence}

Fiche de poste : \textit{``Développeur Python senior, 5~ans, Paris ou remote,
FastAPI, LangGraph, Docker, RAG''} (21 mai 2026).

\begin{table}[H]
\centering
\caption{Métriques du run de référence}
\label{tab:metriques_run}
\begin{tabularx}{\textwidth}{L{5.5cm}C{3cm}L{4.5cm}}
\toprule
\textbf{Métrique} & \textbf{Valeur} & \textbf{Commentaire} \\
\midrule
Profils bruts collectés & 46 & DDG + GitHub + SO \\
Après filtrage A3c & 13 & Bruit : 68\,\% \\
Après déduplication A6 & 12 & 1 doublon nom \\
Candidats évalués (A4) & 10 & Limite MAX\_PROFILS=10 \\
Candidats valides (A5) & 5 & Dont 3 upgrades LinkedIn \\
Messages générés (A7) & 2 & Mode absolu (score $\geq 75$) \\
Durée totale & 173~s & 10 appels LLM parallèles \\
\bottomrule
\end{tabularx}
\end{table}

\begin{table}[H]
\centering
\caption{Meilleurs candidats identifiés}
\label{tab:candidats}
\begin{tabularx}{\textwidth}{L{4.5cm}C{1.5cm}C{2cm}L{5cm}}
\toprule
\textbf{Profil} & \textbf{Score} & \textbf{Source} & \textbf{Remarques A5} \\
\midrule
Srikanth Vanjre & 88 & LinkedIn & Senior Backend AI/RAG \\
Shubhangi Gaur & 80 & LinkedIn & Python/FastAPI senior \\
Nikita Koznev & 72 & LinkedIn & Upgrade auto \\
Jyotsna Konchada & 68 & LinkedIn & Upgrade auto \\
Vamshi Krishna & 68 & LinkedIn & Upgrade auto \\
\bottomrule
\end{tabularx}
\end{table}

\section{Apport du RAG sur les runs répétés}

\medskip
\colorbox{lightblue}{\parbox{0.95\textwidth}{
\smallskip
\textbf{\color{primary}Run 1 (RAG vide)} : 0 URL blacklistée, 0 injection, durée 173~s.\\[4pt]
\textbf{\color{primary}Run 2 (RAG peuplé)} : profils SO invalides filtrés avant scraping
(économie 3 appels A4+A5), 5 candidats valides réinjectés.
Gain estimé : $-$6 appels Groq, durée réduite de 30--40\,\%.
\smallskip
}}
\medskip

\section{Limites}

\paragraph{Scraping LinkedIn.} 100\,\% des profils \texttt{/in/} bloqués.
Le guard A5 (upgrade si score $\geq 75$) compense partiellement.

\paragraph{Quota Groq.} 100\,000 tokens/jour $\approx$ 3--5 runs complets.

\paragraph{Calibration RAG.} Seuils calibrés empiriquement sur peu de runs.

\section{Discussion AI4SE / SE4AI}

\paragraph{SE4AI.} Patterns SMA formels, guards déterministes, 91 tests,
métriques par agent, feature model.

\paragraph{AI4SE.} Analyse automatique de fiches (A2), requêtes optimisées (A3a),
scoring sémantique (A4), messages personnalisés (A7).

% =============================================================================
\chapter{Conclusion et perspectives}
\label{chap:conclusion}
% =============================================================================

\section{Bilan}

Ce TER a atteint les trois objectifs :

\begin{enumerate}
  \item \textbf{État de l'art} : 17 articles analysés, convergences SMA/Agentic AI,
        comparaison de 5 frameworks, trajectoire évolutive.
  \item \textbf{Feature model} : variabilité en 5 dimensions, décisions justifiées.
  \item \textbf{Prototype opérationnel} : 8 agents, LangGraph end-to-end (173~s),
        RAG inter-runs (blacklist + cache + injection), HITL, microservices Docker,
        API SSE, UI NiceGUI, 91 tests.
\end{enumerate}

Ce prototype démontre qu'une Agentic AI conçue avec les outils du génie logiciel
(patterns SMA, blackboard, tests, feature model) produit un système fiable, traçable
et évolutif.

\section{Perspectives}

\paragraph{Kubernetes.} Chaque service est sans état local : scaling horizontal
d'A4/A5 directement applicable.

\paragraph{API LinkedIn officielle.} Éliminerait la limitation principale.

\paragraph{Calibration automatique RAG.} Régression logistique sur les décisions
HITL pour calibrer les seuils de similarité.

\paragraph{MARL.} Adaptation des stratégies de requête A3a selon les résultats
passés \citep{cmupaper2021}.

\paragraph{Gouvernance.} Registre de décisions pour traçabilité et explicabilité
\citep{milosevic2026}.

% ─────────────────────────────────────────────────────────────────────────────
\bibliographystyle{apalike}
\bibliography{biblio}
% ─────────────────────────────────────────────────────────────────────────────

\end{document}
